\documentclass[12pt, a4paper, oneside]{report}
% ============================================================
%  Shared preamble — ViKi BS thesis (Innopolis University)
%  Compile with: xelatex (twice) — Cyrillic requires XeLaTeX
% ============================================================

\usepackage{fontspec}
\usepackage{polyglossia}
\setmainlanguage{english}
\setotherlanguage{russian}

\setmainfont{FreeSerif}[Ligatures=TeX]
\setsansfont{FreeSans}[Ligatures=TeX]
\setmonofont{DejaVu Sans Mono}[Scale=0.85]
\newfontfamily\cyrillicfont{FreeSerif}[Ligatures=TeX]
\newfontfamily\cyrillicfontsf{FreeSans}[Ligatures=TeX]
\newfontfamily\cyrillicfonttt{DejaVu Sans Mono}[Scale=0.85]

\usepackage{geometry}
\geometry{a4paper, top=25mm, bottom=25mm, left=30mm, right=15mm}

\usepackage{amsmath, amssymb, amsfonts, bm}
\usepackage{graphicx}
\usepackage{booktabs}
\usepackage{longtable}
\usepackage{array}
\usepackage{multirow}
\usepackage{caption}
\usepackage{subcaption}
\usepackage{enumitem}
\usepackage{xcolor}
\usepackage{setspace}
\usepackage{tikz}
\usetikzlibrary{arrows.meta, positioning, shapes.geometric, fit, backgrounds}
\usepackage[hidelinks]{hyperref}
\usepackage{tocloft}

\onehalfspacing
\setlength{\parindent}{1.25em}
\setlength{\parskip}{0.25em}

% --- IEEE-ish caption style -------------------------------
\captionsetup[figure]{labelfont=bf, labelsep=period, name=Fig., font=small}
\captionsetup[table]{labelfont=bf, labelsep=newline, justification=centering,
                     name=TABLE, font=small}
\renewcommand{\thetable}{\Roman{table}}

% --- placeholder macro ------------------------------------
% Every unresolved number / empirical claim is wrapped in \ph{}
% so it can be found with a single grep before submission.
\definecolor{phcol}{RGB}{176,0,32}
\newcommand{\ph}[1]{\textcolor{phcol}{\textbf{#1}}}
\newcommand{\TODO}[1]{\textcolor{phcol}{\small\textsf{[TODO: #1]}}}

% --- figure placeholder box -------------------------------
\newcommand{\figplaceholder}[2][0.62\textwidth]{%
  \begin{center}
  \fbox{\parbox[c][#2][c]{#1}{\centering\textcolor{phcol}{\textsf{\small
  [placeholder — insert figure]}}}}
  \end{center}}

% --- math notation ----------------------------------------
\newcommand{\R}{\mathbb{R}}
\newcommand{\SE}{\mathrm{SE}(3)}
\newcommand{\so}{\mathfrak{se}(3)}
\newcommand{\bq}{\bm{q}}
\newcommand{\bT}{\bm{T}}
\newcommand{\be}{\bm{e}}
\newcommand{\bx}{\bm{x}}
\newcommand{\bw}{\bm{w}}
\newcommand{\bJ}{\bm{J}}
\newcommand{\norm}[1]{\left\lVert #1 \right\rVert}
\DeclareMathOperator{\Log}{Log}
\DeclareMathOperator*{\argmin}{arg\,min}

% --- misc --------------------------------------------------
\newcommand{\viki}{ViKi}
\renewcommand{\cftsecleader}{\cftdotfill{\cftdotsep}}


\setmainlanguage{russian}
\setotherlanguage{english}

\begin{document}

\begin{titlepage}
\centering
\vspace*{10mm}
{\large УНИВЕРСИТЕТ ИННОПОЛИС}\\[2mm]
{\normalsize Бакалавриат по направлению «Компьютерные науки»}\\[35mm]

{\LARGE\bfseries ViKi: кинематическая имитация\\[3mm]
на основе компьютерного зрения\par}
\vspace{6mm}
{\large Краткое изложение --- главы 1--3\par}
\vspace{20mm}

\begin{tabular}{rl}
\textit{Автор:} & \ph{Артём [Фамилия]}\\[4mm]
\textit{Научный руководитель:} & \ph{[ФИО]}\\[4mm]
\textit{Группа:} & \ph{[группа]}\\
\end{tabular}

\vfill
{\normalsize Иннополис, Россия\\ \the\year}
\end{titlepage}

\tableofcontents
\cleardoublepage

% =========================================================
\chapter{Постановка задачи и актуальность}

\section{Предметная область}

Роботов учат новым задачам манипуляции по демонстрациям человека
(imitation learning). Чем больше и чище демонстрации, тем лучше работает
итоговая политика, например Action Chunking Transformer (ACT) или Diffusion
Policy.

Сегодня демонстрации чаще всего собирают \emph{телеоперацией}: человек ведёт
робота по траектории задачи, а позы суставов и видео записываются. Плюс в том,
что такие данные сразу исполнимы роботом. Минусы: на каждый эпизод уходит
реальное время оператора; нужно дорогое и специфичное оборудование (VR-контроллер,
экзоскелет, устройство с силовой отдачей); записанное движение определяется не
столько задачей, сколько задержками и эргономикой интерфейса.

\emph{Пассивная запись} обращает этот компромисс: человек выполняет задачу голыми
руками, без интерфейса. Это дёшево, быстро и естественно --- за время одного
телеоперационного эпизода демонстратор успевает записать десятки. Но теряется как
раз то, что телеоперация гарантирует: нет меток действий, а кинематика руки
человека не совпадает с кинематикой робота. Устранение этого несовпадения ---
задача \emph{ретаргетинга} --- и есть тема работы.

\section{Проблема}

Пассивный захват даёт наблюдения, а не действия. Нужно превратить траекторию
кисти человека в траекторию суставов робота, которая одновременно: (i) повторяет
значимое для задачи движение, (ii) достаточно гладкая для обучения политикой и
(iii) физически исполнима на конкретном манипуляторе. Задача недоопределена и
осложнена шумом восприятия. Преобладают два режима отказа.

\textbf{Шум усиливается через кинематическую цепь.} Оценка глубины ключевых точек
кисти шумит на масштабе сантиметров. Обратная кинематика нелинейна, поэтому малое
возмущение в пространстве задачи вблизи сингулярности отображается в большое
отклонение в пространстве суставов. Привычный ответ --- сгладить наблюдаемую
траекторию до решения обратной кинематики и сгладить траекторию суставов после ---
не помогает: сглаживание применяется в пространстве задачи, а гладкость нужна в
пространстве суставов, и вдобавок вместе с шумом удаляется истинное быстрое
движение.

\textbf{Зависимость от представления.} Если хранить траекторию как абсолютные
позиции в системе координат камеры или мира, она привязана к раскладке конкретного
сеанса записи. Сдвинули объект --- траектория неверна; сменили робота --- неверна
иначе. Обобщается отношение «кисть--объект», а не абсолютный путь кисти в
пространстве.

\section{Научный пробел}

У обоих режимов отказа есть известные решения, но они приняты в разных, почти не
пересекающихся сообществах.

\emph{Оптимизационный ретаргетинг} --- когда точность соответствия данным,
гладкость во времени и ограничения на суставы ставятся как одна задача условной
оптимизации по конфигурации робота --- стал стандартом в системах
\emph{онлайн-телеоперации}. Каноническую форму задали AnyTeleop и библиотека
\texttt{dex-retargeting}; на них опираются Open-TeleVision, Bunny-VisionPro, ACE;
DexFlow расширяет штраф гладкости до второго порядка. Но человек всегда остаётся в
контуре, оптимизация идёт на интерактивных частотах, и датасет не формируется.

\emph{Офлайн-пайплайны}, делающие датасеты из пассивного видео человека, берут
объектно-центричные представления, но не единую оптимизационную постановку. ORION
планирует по разреженным ключевым кадрам объекта и покадровый ретаргетинг кисти не
выполняет вовсе. OKAMI восстанавливает модель тела и кисти и ретаргетит в две
развязанные стадии --- варпинг референсной траектории под новое положение объекта,
затем решение обратной кинематики, --- то есть ту самую последовательную схему,
которую критикуют выше.

Целевой анализ литературы не нашёл системы, объединяющей сразу: пассивный
безмаркерный захват; метрически откалиброванный мультикамерный RGB-D риг;
объектно-центричное плотное траекторное представление; ретаргетинг как единый
функционал стоимости; экспорт готового к обучению датасета; и физическое
воспроизведение как этап валидации до экспорта. Ближайшие системы закрывают часть
критериев и не закрывают остальные. Пробел, таким образом, не в новом операторе
ретаргетинга, а в переносе устоявшейся формулировки в режим, где её не применяли,
вместе с нужной этому режиму сенсорной и валидационной инфраструктурой.

\section{Исследовательские вопросы и гипотезы}

\noindent\textbf{Основной вопрос.} Как демонстрации манипуляции, захваченные
пассивно откалиброванным мультикамерным RGB-D ригом, ретаргетить на робота так,
чтобы получать датасеты для обучения визуомоторных имитационных политик?

\begin{enumerate}[label=\textbf{ИВ\arabic*.}, leftmargin=2.6em]
  \item Как единый взвешенный функционал стоимости (член соответствия данным со
        взвешиванием по достоверности плюс штрафы на скорость и ускорение) влияет
        на точность и гладкость траекторий суставов по сравнению с последовательной
        схемой smooth--IK--smooth?
  \item Какая точность позы запястья достижима при взвешенном по достоверности
        слиянии нескольких откалиброванных RGB-D видов с общим началом координат по
        доске ChArUco, и как раскладывается остаточная ошибка?
  \item Насколько воспроизведение синтезированных траекторий на реальном
        манипуляторе до экспорта повышает исполнимость датасета?
\end{enumerate}

\noindent\textbf{Гипотезы.}
\begin{itemize}[leftmargin=2.6em]
  \item \textbf{Г1.} Единый взвешенный функционал даёт меньшую ошибку слежения в
        пространстве задачи и большую гладкость суставов, чем схема
        smooth--IK--smooth, на тех же демонстрациях.
  \item \textbf{Г2.} Взвешивание члена данных по покадровой достоверности
        восприятия снижает ошибку слежения относительно равномерного взвешивания,
        сильнее --- на кадрах с частичной окклюзией кисти.
  \item \textbf{Г3.} Отбор траекторий по результатам воспроизведения до экспорта
        повышает долю эпизодов, исполнимых на роботе.
\end{itemize}

Вопросы проверены по критериям FINER: \emph{выполнимы} (оборудование, решатель и
протокол доступны, класс задач намеренно узкий); \emph{интересны} и
\emph{релевантны} (стоимость демонстраций --- признанное узкое место);
\emph{новы} в смысле переноса режима, а не нового оператора; \emph{этичны} ---
испытуемых, кроме автора, нет, персональные данные не сохраняются.

% =========================================================
\chapter{Обзор литературы}

\section{Как представляют кисть человека}

Часть работ восстанавливает полную кинематику демонстратора: OKAMI (SMPL-H для
тела и MANO для кисти), DexCanvas (подгонка MANO к данным 22 ИК-камер захвата
движения и двух RGB-D сенсоров), DexMV (MANO вместе с позой объекта). Полная
реконструкция богата и нужна для многопальцевых антропоморфных кистей, но тянет
тяжёлый перцептивный стек и добавляет плохо измеримую ошибку --- подгонку меша.
Для параллельного схвата большая часть этого богатства всё равно отбрасывается.

Здесь кисть представлена минимально: поза запястья в $\SE$ плюс бинарное состояние
захвата (открыт/закрыт). Так делает магистральное направление телеоперации ---
Open-TeleVision, Bunny-VisionPro, ACE. Абляции Xin и соавт. (полный функционал
ретаргетинга, из которого члены удаляют по одному) показывают, что точное слежение
за глобальной позой запястья часто избыточно и даже вредит точности по кончикам
пальцев. Отсюда: цель по запястью должна входить в оптимизацию как взвешенный
мягкий член, а не жёсткое ограничение.

\section{Формулировки ретаргетинга}

Преобладающая современная формулировка ставит ретаргетинг как условную задачу
(нелинейный МНК или квадратичное программирование) по конфигурации робота на
каждом кадре: минимум взвешенной суммы члена соответствия данным и члена
гладкости при ограничениях на пределы суставов. Каноническая реализация ---
AnyTeleop.

Два уточнения важны для настоящей работы. Первое --- покадровое взвешивание
членов: в Open-TeleVision невязка каждой ключевой точки масштабируется весом,
зависящим от расстояния; это прямой аналог взвешивания по достоверности
восприятия. Второе --- гладкость второго порядка: DexFlow штрафует кривизну, а не
только скорость. Робастная потеря Хьюбера на члене данных отбраковывает выбросы
восприятия внутри оптимизатора, а не сглаживает их заранее.

Последовательная схема --- сгладить наблюдаемую траекторию, покадрово решить
обратную кинематику, сгладить выход в пространстве суставов --- распространена в
офлайне. Её дефект структурный: сглаживание идёт в пространстве задачи, а
гладкость нужна в пространстве суставов; из-за нелинейности прямой кинематики
гладкий вход не даёт гладкого выхода, особенно у сингулярностей. В едином
функционале штраф гладкости действует прямо на конфигурацию.

\section{Представление задачи}

Объектно-центричные представления доминируют в недавних работах про обобщение
между сценами и роботами. ORION задаёт задачу как 3D-траекторию объекта
относительно его начального и конечного положений и утверждает, что это устойчиво
к смене фона, сдвигу камеры и новым экземплярам объекта. SPOT кодирует задачу как
$\SE$-траекторию позы объекта относительно цели, чтобы отвязать действия от
сенсорного входа. Подходы на основе потока дают тот же аргумент без предположения
о жёсткости объекта.

Компромисс: объектно-центричные представления требуют устойчивой оценки позы или
сегментации объекта, а траектории в мировой системе координат проще, но привязаны
к раскладке съёмки. ViKi хранит оба представления --- выбор остаётся за
потребителем данных, а стоимость хранения мала.

\section{Калибровка, слияние и синхронизация RGB-D рига}

Доска ChArUco сочетает субпиксельную точность углов шахматной доски с устойчивым
распознаванием маркеров ArUco. Доска $6\times8$ даёт до 35 углов на вид, остаётся
детектируемой при частичной окклюзии и на краю кадра; типичная ошибка репроекции
--- ниже половины пикселя.

Romeo и соавт. дают чистое сравнение процедур калибровки мульти-Kinect ригов:
трёхмерные процедуры точнее двумерных на основе шахматной доски (расхождения
облаков точек порядка $10$--$21$~мм). Bu и Lee отмечают, что уточнение методом
итеративной ближайшей точки не сходится для почти противоположных ракурсов из-за
малого перекрытия облаков.

Сами сенсоры ограничивают точность снизу: Azure Kinect DK специфицирован
случайной ошибкой не более $17$~мм и систематической ниже $11$~мм плюс
$0{,}1$~\% дистанции. Практический вывод: между откалиброванными видами следует
ожидать согласования сантиметрового, а не миллиметрового уровня.

Перекрывающиеся измерения глубины лучше объединять не наивным средним, а
взвешенным по достоверности: вклад тем больше, чем ближе точка и чем ближе к
перпендикуляру наблюдается поверхность.

Часы камер разных семейств несопоставимы. Документированная практика --- метить
каждый входящий кадр единым монотонным таймстампом хоста и выравнивать потоки по
ближайшему соседу в этом общем домене (точность миллисекундного масштаба).
Триггеры Azure Kinect и RealSense несовместимы, поэтому другого механизма
выравнивания двух семейств нет.

\section{Место предлагаемой системы}

Сравнение --- по шести критериям, которые вместе задают архитектуру: (К1)
пассивный безмаркерный захват без телеоперационного оборудования; (К2)
мультикамерное метрически откалиброванное RGB-D зондирование; (К3)
объектно-центричное плотное траекторное представление; (К4) ретаргетинг как
единый взвешенный по достоверности функционал стоимости; (К5) экспорт датасета для
имитационного обучения; (К6) физическое воспроизведение для валидации до экспорта.

Ни одна найденная система не закрывает все шесть. HOWTransfer ближе всего по
порядку операций (К1, К6), но кисте-центрична, использует стереопару RGB вместо
метрического RGB-D и не применяет единый функционал. Vi-RoIL ближе по
представлению (К3), но восстанавливает масштаб из монокулярного вида маркерного
куба и проверяется только в симуляции. Пересечение онлайн-оптимизационного
ретаргетинга и офлайн-генерации объектно-центричных датасетов пустует --- на него
и нацелена работа.

% =========================================================
\chapter{Предлагаемое решение}

\section{Структура пайплайна}

ViKi --- офлайн-пайплайн: запись, обработка и экспорт идут отдельными стадиями, и
ни одна не обязана работать в реальном времени. Это осознанный выбор: без
ограничения реального времени можно оптимизировать сразу по всей траектории,
итеративно уточнять и добавить стадию проверки на реальном оборудовании.

Стадий шесть:
\begin{enumerate}[leftmargin=2.2em]
  \item захват синхронизированных RGB-D потоков;
  \item восприятие --- детекция ключевых точек кисти, подъём в 3D по данным
        глубины, слияние видов со взвешиванием по достоверности;
  \item представление траектории запястья относительно объекта и относительно
        якоря рабочей области;
  \item ретаргетинг --- минимизация единого взвешенного функционала при
        кинематических и коллизионных ограничениях;
  \item валидация воспроизведением на физическом манипуляторе;
  \item экспорт в формате, совместимом с LeRobot.
\end{enumerate}

\section{Калибровка и синхронизация}

Внутренние параметры всех цветовых и глубинных потоков оцениваются по наблюдениям
доски ChArUco. Внешняя калибровка следует принципу \emph{якоря рабочей области}:
каждая камера калибруется относительно доски ChArUco, жёстко закреплённой в
рабочей области и задающей общую мировую систему координат, --- а не относительно
опорной камеры. Следствия: добавление или перемещение камеры не аннулирует
калибровку остальных, а базу манипулятора можно привязать к тем же координатам без
отдельной цепочки камера--камера.

Две камеры Azure Kinect синхронизируются аппаратно в режиме ведущий/ведомый:
ведомые запускаются раньше ведущего, а экспозиции глубины разнесены во времени,
чтобы лазеры не засвечивали друг друга. RealSense этот триггер не поддерживает,
поэтому между семействами используется единый монотонный таймстамп хоста на каждый
кадр с последующим сопоставлением по ближайшему соседу. Часы берутся монотонные, а
не реального времени: последние может скачком поправить служба синхронизации
времени и испортить вычисление интервалов посреди записи.

\section{Слияние со взвешиванием по достоверности}

Оценки от нескольких видов объединяются взвешенным средним. Вес каждой ключевой
точки в каждом виде --- произведение четырёх множителей: оценки видимости
детектора; валидности глубины сенсора в соответствующем пикселе; обратного
квадрата дальности; косинуса угла падения, отсечённого снизу нулём. Форма выбрана
так, чтобы нуль в любом множителе --- точка не найдена, пиксель глубины невалиден,
скользящий ракурс --- полностью убирал вклад этого вида, а не просто ослаблял его.

\section{Ретаргетинг как единый функционал стоимости}

Пусть $\bq_t \in \R^{n}$ --- конфигурация манипулятора в момент $t$, а
$f(\bq_t) \in \SE$ --- прямая кинематика схвата. Ошибка в пространстве задачи
относительно желаемой позы --- логарифм ошибки позы,
$\be_t(\bq_t) = \Log\big((\bT^{R,\mathrm{des}}_{E,t})^{-1} f(\bq_t)\big)^{\vee}
\in \R^{6}$. Ретаргетинг ставится как одна условная задача:

\begin{equation}
  \min_{\bq_{1:T}} \sum_{t=1}^{T} \Big[
    \omega_t \rho_\delta\!\left(\norm{\be_t(\bq_t)}_{\bm{\Sigma}}\right)
    + \lambda_v \norm{\bq_t - \bq_{t-1}}^{2}
    + \lambda_a \norm{\bq_t - 2\bq_{t-1} + \bq_{t-2}}^{2}
    + \lambda_r \norm{\bq_t - \bq_{\mathrm{ref}}}^{2}
  \Big]
\end{equation}

при ограничениях на пределы суставов и скоростей и с барьерными функциями
коллизий. Здесь $\omega_t$ --- покадровый вес достоверности восприятия,
$\rho_\delta$ --- потеря Хьюбера, $\lambda_v, \lambda_a, \lambda_r$ --- веса
штрафов на скорость, ускорение и отклонение от опорной позы.

Вес $\omega_t$ собирается из достоверностей ключевых точек, породивших оценку
запястья в момент $t$. На плохо наблюдаемых кадрах он мал, и траекторию через них
проводят члены гладкости, а не член данных, стягивающий её к ложному наблюдению.
Потеря Хьюбера ограничивает влияние грубых выбросов прямо в члене данных --- это и
заменяет предварительное сглаживание.

Ключевое отличие от последовательной схемы: штрафы гладкости действуют на
конфигурацию, то есть в том пространстве, где гладкость и нужна, и балансируются
против соответствия данным самим решателем. На кадрах высокой достоверности вес
велик и решение отслеживает наблюдение; на кадрах низкой доминируют члены
гладкости и решение интерполирует.

Задача решается средствами Pink на Pinocchio: каждый член --- взвешенная
QP-задача, ограничения входят как неравенства, избегание коллизий --- через
барьерные функции, а демпфирование Левенберга--Марквардта стабилизирует решение у
сингулярностей. Так как режим офлайновый, задача решается по всей траектории, а не
покадрово: член ускорения корректно определён на каждом внутреннем шаге, а решение
не стеснено требованием причинности.

Динамические фильтры, налагающие модель движения (фильтр Калмана и подобные),
намеренно не применяются: они вносят априорное предположение о динамике, которое
затем переоценивает оптимизатор, а взаимодействие двух априоров трудно
охарактеризовать.

\section{Валидация воспроизведением}

Синтезированные траектории исполняются на физическом манипуляторе до финализации
датасета. Стадия служит трём целям:

\begin{enumerate}[leftmargin=2.2em]
  \item \emph{Отсев неисполнимого.} Траектория может удовлетворять кинематическим
        ограничениям в модели и всё же отказать на оборудовании --- из-за пределов
        отслеживания контроллера, немоделированных коллизий или границы рабочей
        области.
  \item \emph{Закрытие разрыва эмбодиментов у источника.} Записываются
        проприоцептивные состояния, которых робот действительно достиг, а не
        предсказанные оптимизатором. Экспорт именно записанных состояний даёт
        честные пары «наблюдение--действие» с целевого робота.
  \item \emph{Сигнал для курации.} Эпизоды ранжируются по невязке отслеживания при
        воспроизведении.
\end{enumerate}

\end{document}
